\documentclass[11pt]{article}
% Shared preamble for the MSH_Tensor_Network_Methods function references (compile with tectonic / XeLaTeX)
\usepackage{fontspec}
\usepackage{unicode-math}
\setmainfont{texgyrepagella}[Extension=.otf,UprightFont=*-regular,BoldFont=*-bold,ItalicFont=*-italic,BoldItalicFont=*-bolditalic]
\setmathfont{texgyrepagella-math.otf}
\setmonofont{DejaVu Sans Mono}[Scale=0.82]
\newfontfamily\symfb{DejaVu Sans}[Scale=0.82]
\usepackage{newunicodechar}
\newunicodechar{ℓ}{{\symfb ℓ}}
\newunicodechar{∥}{{\symfb ∥}}
\newunicodechar{≪}{{\symfb ≪}}
\usepackage[a4paper,margin=2.2cm]{geometry}
\usepackage{amsmath}
\usepackage{booktabs,tabularx,longtable,array}
\usepackage{graphicx}
\usepackage{xcolor}
\usepackage{float}
\usepackage{caption}
\captionsetup{font=small,labelfont=bf}
\usepackage{fvextra}
\fvset{breaklines,breakanywhere,fontsize=\small,frame=leftline,framerule=1.5pt,rulecolor=\color{gray!45},xleftmargin=6pt}
\usepackage{enumitem}
\setlist{nosep,leftmargin=*}
\usepackage[hidelinks]{hyperref}
\usepackage{titlesec}
\titleformat{\section}{\Large\bfseries}{\thesection}{0.6em}{}
\titleformat{\subsection}{\large\bfseries}{\thesubsection}{0.6em}{}
\titleformat{\subsubsection}{\normalsize\bfseries}{\thesubsubsection}{0.6em}{}
\setlength{\parskip}{0.45em}\setlength{\parindent}{0pt}
% one box per function: \fn{name}{signature} then a description
\newcommand{\fnname}[1]{\texttt{\textbf{#1}}}
\newenvironment{function}[1]{\par\medskip\noindent\textbf{\large\texttt{#1}}\par\nopagebreak}{\par\medskip}
\newcommand{\code}[1]{\texttt{#1}}

\graphicspath{{figs/}}
\setlength{\emergencystretch}{3em}
\newcommand{\nb}{\code{kitaev\_charge\_difference.ipynb}}
\newcommand{\ket}[1]{\lvert #1\rangle}
\newcommand{\bra}[1]{\langle #1\rvert}
\newcommand{\Tr}{\operatorname{Tr}}

\begin{document}

\begin{titlepage}
\centering
\vspace*{3cm}
{\Huge\bfseries Kitaev Chain Charge Difference\par}
\vspace{0.8cm}
{\LARGE Function reference\par}
\vspace{1.2cm}
{\large Notebook: \nb\par}
\vspace{0.3cm}
{\large Repository: \code{MSH\_Tensor\_Network\_Methods}\par}
\vspace{0.3cm}
{\large Julia 1.13, TensorBinding.jl, ITensors.jl / ITensorMPS.jl\par}
\vspace{1.5cm}
{\large 2026-10-29\par}
\vfill
\begin{minipage}{0.85\textwidth}\small
This document explains every function and helper defined in the notebook: what it computes, the method,
inputs and outputs, the numerical settings that matter, and where it is used. All statements are taken from the
notebook's code, comments, markdown and saved outputs. Where a library routine of TensorBinding.jl is relevant, its
behaviour is quoted from the package source and marked as such.
\end{minipage}
\end{titlepage}

\tableofcontents
\clearpage

%======================================================================
\section{Overview}
%======================================================================

\subsection{The model as built in the notebook}\label{sec:model}

The notebook studies a spinless one-dimensional Kitaev chain with nearest-neighbour hopping, a site-dependent
on-site potential $V(x)$ (called the chemical potential in the notebook) and nearest-neighbour $p$-wave pairing.
Every Hamiltonian is created with three TensorBinding calls:
\begin{Verbatim}
H = TensorBinding.get_Hamiltonian("chain_1d", 1.0; L = Lℓ)   # chain of 2^L sites, hopping 1.0
TensorBinding.add_onsite!(H, Vf)                               # V(x); Vf receives a 0-indexed site
TensorBinding.add_superconductivity!(H, 1.0)                   # pairing amplitude 1.0
\end{Verbatim}
For spinless fermions \code{add\_superconductivity!} prints \emph{``On-site (s-wave) pairing is forbidden for
spinless fermions \ldots{} Constructing nearest-neighbour p-wave instead''}, so the pairing is nearest-neighbour
$p$-wave with $\Delta = 1$. The resulting object is a \code{TBHamiltonian} (printed as
\code{L=6, N=64 +BdG, maxlinkdim=5}) whose field \code{.mpo} is the Bogoliubov--de~Gennes (BdG) Hamiltonian as an
MPO, \code{.sites} its site indices, \code{.L} the number of position bits and \code{.N}$=2^L$ the number of lattice
sites.

\paragraph{Sign conventions of the library Hamiltonian.}
Cell 35 reads the dense BdG matrix of the hopping-plus-pairing part off the MPO (\code{dense\_bdg}, \S\ref{sec:td})
and compares it element by element with
\[
  H_{\text{BdG}}=\begin{pmatrix} h & \Delta\\ \Delta^\dagger & -h^{T}\end{pmatrix},\qquad
  h_{x,x\pm1}=+t,\quad \Delta_{x,x+1}=+\Delta,\quad \Delta_{x+1,x}=-\Delta,\quad \Delta_{x,x}=0,
\]
with $t=\Delta=1$ (maximum deviation $3.6\times10^{-15}$). In second quantization,
\[
  H=\sum_x\Bigl[t\,(c_x^\dagger c_{x+1}+\text{h.c.})+\Delta\,(c_x^\dagger c_{x+1}^\dagger+\text{h.c.})+V(x)\,c_x^\dagger c_x\Bigr],
\]
with the on-site term entering the BdG matrix as $\tau_z\otimes\operatorname{diag}V$. The hopping thus enters with a
$+$ sign. The gauge transformation $c_x\to(-1)^x c_x$ maps $t\to-t$ (and $\Delta\to-\Delta$), so the spectrum is that of
the usual Kitaev chain with hopping $-t$, topological for $|V|<2|t|$.

\paragraph{Nambu and position qubits.}
A BdG single-particle state is an MPS on $L+1$ two-level sites. Site~1 is the Nambu index
($\ket{0}$ = particle, $\ket{1}$ = hole; ``Nambu is first index'' in the code). Sites $2,\dots,L+1$ are the binary
digits $b_1\dots b_L$ of the position, most significant first, with
\[
  x = 1 + \sum_{k=1}^{L} b_k\, 2^{L-k}, \qquad x = 1,\dots,N=2^L .
\]
A BdG vector $\psi = (u, v)$ thus has particle amplitudes $u(x) = \langle 0, x\vert\psi\rangle$ and hole amplitudes
$v(x) = \langle 1, x\vert \psi\rangle$. The on-site potential enters as $\tau_z\otimes\operatorname{diag}V(x)$; the
notebook checks this numerically ($\langle 0,x\vert H\vert 0,x\rangle = V(x)$, see \code{charge\_difference} and the
time-dependent section).

\paragraph{Particle--hole conjugation.}
The BdG Hamiltonian anticommutes with $C=\tau_x K$ ($K$ = complex conjugation). On an MPS the notebook implements
$C$ as \code{dag(ψ)} (complex conjugation of all tensors) followed by $\tau_x$ applied to the Nambu site. For
$\psi=(u,v)$ this gives $C\psi = (v^*, u^*)$. If $H\psi = E\psi$ then $H C\psi = -E\, C\psi$, so states come in $\pm E$
pairs. A Majorana state is self-conjugate, $C\gamma=\gamma$.

\paragraph{Parameters.}
\begin{center}
\begin{tabular}{@{}lll@{}}
\toprule
Quantity & Static part (\S\ref{sec:fermion}--\ref{sec:dq}) & Time-dependent part (\S\ref{sec:td}) \\
\midrule
Position bits $L$ & \code{Lℓ = 6} & \code{Lt = 5} \\
Sites $N = 2^L$ & 64 & 32 \\
Hopping, pairing & 1.0, 1.0 & 1.0, 1.0 \\
Topological potential & \code{V\_in = -1.0} & \code{μ\_topo\_t = -0.4} \\
Trivial potential & \code{V\_out = -4.0} & \code{μ\_triv\_t = -5.6} \\
Topological region & centred, length $\ell$ (\code{region(ℓ)}) & whole chain at $t=0$, site 16 at $t=T$ \\
\bottomrule
\end{tabular}
\end{center}
With hopping 1 the notebook's markdown refers to the topological criterion $|\mu|<2t$; $V=-1$ and $-0.4$ are
inside it, $-4$ and $-5.6$ outside.

\subsection{Logical states and the charge difference}

Following the notebook's markdown, the two logical states $\ket{0}$ (fermion mode empty) and $\ket{1}$ (mode
occupied) are represented by their BdG correlation matrices
\[
  \Gamma_{ab} = \langle \Psi_b^\dagger \Psi_a\rangle, \qquad
  \Psi = (c_1,\dots,c_N, c_1^\dagger,\dots,c_N^\dagger).
\]
Since $H$ is quadratic, $\Gamma$ fully characterises the BCS state and is the projector on its occupied BdG states.
Both logical states share the filled bulk projector $P_{\text{bulk}}$ and differ in which member of the near-zero
pair $\pm E$ is filled:
\[
  \Gamma_0 = P_{\text{bulk}} + \ket{\bar\psi_f}\bra{\bar\psi_f}, \qquad
  \Gamma_1 = P_{\text{bulk}} + \ket{\psi_f}\bra{\psi_f}, \qquad \bar\psi_f = C\psi_f .
\]
The local density is the particle--particle diagonal, $n(x)=\langle c_x^\dagger c_x\rangle = \langle 0,x\vert\Gamma\vert 0,x\rangle$.
Because $\Gamma_1-\Gamma_0 = \ket{\psi_f}\bra{\psi_f} - \ket{\bar\psi_f}\bra{\bar\psi_f}$ exactly,
\begin{equation}\label{eq:dq}
  \delta n(x) = n_1(x)-n_0(x) = |u_f(x)|^2 - |u_{\bar f}(x)|^2 = |u_f(x)|^2-|v_f(x)|^2,
  \qquad
  \Delta Q = \sum_x \delta n(x) = \bra{\psi_f}\tau_z\ket{\psi_f}.
\end{equation}
The second equality uses $u_{\bar f} = v_f^*$ from $C\psi=(v^*,u^*)$. The physical charge is $-e\,\delta n$;
plots label $|\Delta Q|$ in units of $-e$.

\subsection{Workflow of the notebook}
\begin{enumerate}
\item Helpers to build position basis states (\S\ref{sec:basis}).
\item ``Good'' single-particle states: penalized DMRG for the near-zero doublet, localisation of the Majoranas by
      diagonalising the position operator, self-conjugation, pairing into $\psi_f$ (\S\ref{sec:fermion}).
\item KPM local density of states and the bulk gap estimate \code{kpm\_gap} (\S\ref{sec:kpm}).
\item McWeeny-purified bulk projector, $\Gamma_0$, $\Gamma_1$, consistency checks and local densities (\S\ref{sec:gamma}).
\item $\Delta Q$ versus the length $\ell$ of the topological region, by two methods (\S\ref{sec:dq}).
\item Time-dependent smoothstep ramp that moves the Majoranas together: TDVP evolution of $\psi_f$, $\Delta Q(t)$,
      $\delta n(x,t)$, and full densities with a TDVP-evolved or McWeeny (adiabatic) bulk (\S\ref{sec:td}).
\end{enumerate}

%======================================================================
\section{Position basis states}\label{sec:basis}
%======================================================================

\begin{function}{to\_bits(n, L)}
\begin{Verbatim}
function to_bits(n::Integer, L::Integer)
    return [parse(Int, c) for c in lpad(string(n; base=2), L, '0')]
end
\end{Verbatim}
Returns the $L$-digit binary representation of $n$ as a \code{Vector\{Int\}}, most significant bit first.
\code{string(n; base=2)} gives the binary string, \code{lpad} pads it with \code{'0'} on the left to length $L$,
and each character is parsed to an integer. Callers pass \code{x - 1} to convert the 1-indexed site $x$ to
$0,\dots,N-1$.
\end{function}

\begin{function}{basis\_state(sites, bits)}
\begin{Verbatim}
function basis_state(sites, bits::Vector{Int})
    return MPS([onehot(sites[i] => bits[i] + 1) for i in eachindex(sites)])
end
\end{Verbatim}
Builds the product MPS $\ket{b_0 b_1\dots b_L}$. \code{onehot(s => 1)} is $[1,0]=\ket{0}$ and
\code{onehot(s => 2)} is $[0,1]=\ket{1}$; the \code{+1} converts bits to Julia's 1-based indices. The bond
dimension is 1. Throughout the notebook it is called as
\code{basis\_state(sites, vcat([nb], to\_bits(x - 1, L)))}, where the leading bit \code{nb} selects the Nambu
sector (0 = particle, 1 = hole). It is the building block of \code{site\_weights}, \code{pbasis}, \code{pb},
\code{pbt}, \code{hbt} and \code{bdg\_basis}.
\end{function}

%======================================================================
\section{The fermion mode from the Majorana doublet}\label{sec:fermion}
%======================================================================

The markdown lists the procedure for ``good'' single-particle states:
(1) find orthonormal $\{\psi_1,\psi_2\}$ spanning the near-zero subspace with DMRG;
(2) localise the Majorana wavefunctions;
(3) enforce self-conjugation $C\gamma=\gamma$;
(4) pair the Majoranas into a fermion.

\begin{function}{doublet(Hmpo, sites; η=1e-2, seed=1234, cutoff=1e-20)}
\begin{Verbatim}
K = TensorBinding.build_K(Hmpo, sites, 0.0, η; cutoff_K=cutoff)
sw = (nsweeps=15, maxdim=[10, 20, 50, 100, 150], cutoff=cutoff,
      noise=[1e-6, 1e-7, 1e-8, 0.0], outputlevel=0)
Random.seed!(seed)
_, ψ1 = dmrg(K, random_mps(sites; linkdims=10); sw...)
_, ψ2 = dmrg(K, [ψ1], random_mps(sites; linkdims=10); weight=10.0, sw...)
ψ1 /= norm(ψ1)
ψ2 = add2(ψ2, -inner(ψ1, ψ2) * ψ1); ψ2 /= norm(ψ2)
\end{Verbatim}
\textbf{Computes} two orthonormal BdG states spanning the near-zero doublet of $H$.

\textbf{Method.} The states with smallest $|E|$ are not the ground state of $H$, so DMRG is run on the positive
operator
\[
  K = (H-\omega)^2 + \eta^2, \qquad \omega=0,
\]
built by \code{TensorBinding.build\_K} (from the package source: $H-\omega\mathbb 1$, squared with \code{apply}
using \code{cutoff\_K} and the default \code{maxdim\_K = 200}, then $+\eta^2\mathbb 1$). Eigenvectors of $K$ are
those of $H$, with eigenvalue $E^2+\eta^2$; the lowest two are the $\pm E$ doublet.

\textbf{Role of $\eta$.} \code{build\_K} serves TensorBinding's spectral DMRG, where $K(\omega,\eta)$ is minimised at a
chosen energy $\omega$ and $\eta$ is a broadening: $\min\langle K\rangle=\eta^2$ when $\omega$ hits an eigenvalue of $H$.
Here only $\omega=0$ is used, and $\eta^2$ shifts every eigenvalue of $K$ by the same constant, so it changes neither the
states DMRG finds nor its convergence (Krylov iterations are invariant under a constant shift). The notebook's comment
records a check on the $\ell=6$ chain: with $\eta=0$, $E_{\text{mode}}$ and $E_{\text{bulk}}$ agree to all printed digits
and the overlap of the two fermion modes is 1. \code{lowest\_bulk} subtracts $\eta^2$ again. The value $\eta=10^{-2}$ is
kept for consistency with the MSH notebooks. \code{cutoff\_K} must stay small (here equal to \code{cutoff}): the package
default $10^{-8}$ can truncate the $\eta^2$ term of the MPO.

\textbf{DMRG runs.} The first DMRG gives $\psi_1$. The second DMRG adds a penalty of
weight 10 against $\psi_1$ (ITensor's \code{dmrg(H, [ψ1], ψ0; weight)}), so it converges to the next eigenvector of
$K$. A final Gram--Schmidt step removes any residual overlap with $\psi_1$.

\textbf{Settings.} 15 sweeps; bond dimension ramped through 10, 20, 50, 100, 150; noise $10^{-6}\to10^{-8}\to0$;
truncation cutoff $10^{-20}$; random initial MPS of bond dimension 10 with a fixed seed (1234) for reproducibility.

\textbf{Returns} \code{(ψ1, ψ2)}. \textbf{Used by} \code{fermion\_mode}, \code{initial\_mode}, \code{final\_branches}.
\end{function}

\begin{function}{lowest\_bulk(Hmpo, sites, ψ1, ψ2; η=1e-2, seed=1234, cutoff=1e-20)}
\begin{Verbatim}
_, φ = dmrg(K, [ψ1, ψ2], random_mps(sites; linkdims=10); weight=10.0, sw...)
φ = add2(φ, -inner(ψ1, φ) * ψ1); φ = add2(φ, -inner(ψ2, φ) * ψ2); φ /= norm(φ)
return φ, sqrt(max(real(inner(φ', K, φ)) - η^2, 0.0))
\end{Verbatim}
\textbf{Computes} the lowest bulk state, i.e.\ the eigenstate with the next-smallest $|E|$ after the doublet (the
bulk band edge), and its energy. It runs a third DMRG on the same $K$ and with the same sweep settings as
\code{doublet}, penalised against both $\psi_1$ and $\psi_2$, then orthogonalises against them. The energy follows
from $\langle K\rangle$:
\[
  |E_{\text{bulk}}| = \sqrt{\max(\bra{\varphi}K\ket{\varphi} - \eta^2,\,0)} .
\]
\textbf{Penalty weight.} The code comment states that the weight (10) must exceed $E_{\text{bulk}}^2+\eta^2$, which
holds for $|E_{\text{bulk}}|<3$; otherwise the penalised doublet states would not lie above the bulk state in the
penalised problem.

\textbf{Returns} \code{(φ, |E\_bulk|)}. \textbf{Used by} \code{fermion\_mode} when \code{gap\_factor} is given.
\end{function}

\begin{function}{add2(a, b)}
\begin{Verbatim}
add2(a, b) = +(a, b; cutoff=1e-16)
\end{Verbatim}
MPS addition with truncation cutoff $10^{-16}$. All linear combinations of MPS in the notebook use it, so that sums
are compressed without losing accuracy.
\end{function}

\begin{function}{fermion\_mode(Hmpo, sites, L; η=1e-2, seed=1234, cutoff=1e-20, gap\_factor=nothing)}
\textbf{Computes} the BdG wavefunction of the ordinary fermion $\psi=(\gamma_L+i\gamma_R)/\sqrt2$ built from the two
Majoranas, its energy, the Majoranas themselves, and optionally the bulk edge.

\textbf{Step 1: doublet.} \code{ψ1, ψ2 = doublet(...)}.

\textbf{Step 2: localisation.} DMRG returns an arbitrary basis of the doublet. The Majoranas are the combinations
that diagonalise the position operator in that subspace. With the bit projector $\hat n_k = \ket{1}\bra{1}$ on
position bit $k$ (MPS site $k+1$),
\[
  \hat x = \mathbb 1 + \sum_{k=1}^{L} 2^{L-k}\, \hat n_k ,
  \qquad
  S_{ab} = \langle\psi_a\vert\psi_b\rangle,\quad
  X_{ab} = \bra{\psi_a}\hat x\ket{\psi_b}, \quad a,b\in\{1,2\}.
\]
\begin{Verbatim}
nk(k) = ITensor([0.0 0.0; 0.0 1.0], sites[k+1]', sites[k+1])
on_site(op, ψ, j) = (φ = copy(ψ); φ[j] = noprime(op * φ[j]); φ)
S = [inner(a, b) for a in ψs, b in ψs]
X = S + sum(2.0^(L - k) * [inner(a, on_site(nk(k), b, k + 1)) for a in ψs, b in ψs] for k in 1:L)
F = eigen(Hermitian((X + X') / 2), Hermitian((S + S')/2))
\end{Verbatim}
The $S$ term supplies the $\mathbb 1$ in $\hat x$. The generalized Hermitian eigenproblem $Xc=\lambda Sc$ is solved
(symmetrised to remove round-off). Eigenvalues are ascending, so column 1 (smaller $\langle x\rangle$) gives
$\gamma_L$ and column 2 gives $\gamma_R$.

\textbf{Step 3: self-conjugation.} Each combination carries an arbitrary phase: $C\gamma = e^{i\theta}\gamma$.
\begin{Verbatim}
τx_s = ITensor([0.0 1.0; 1.0 0.0], sites[1]', sites[1])   # Nambu is first index
C(ψ) = (φ = dag(ψ); φ[1] = noprime(τx_s * φ[1]); φ)
fix(γ) = exp(im * angle(inner(γ, C(γ))) / 2) * γ
\end{Verbatim}
With $\theta=\arg\langle\gamma\vert C\gamma\rangle$, the state $\gamma'=e^{i\theta/2}\gamma$ satisfies
$C\gamma' = e^{-i\theta/2}e^{i\theta}\gamma = \gamma'$.

\textbf{Step 4: pairing and branch choice.}
\[
  \psi = \frac{\gamma_L + i\gamma_R}{\lVert\gamma_L + i\gamma_R\rVert},\qquad E=\bra{\psi}H\ket{\psi},\qquad
  \text{if } E<0:\ \psi\leftarrow C\psi,\ E\leftarrow -E .
\]
The comment notes that $\lVert\gamma_L+i\gamma_R\rVert^2=2$, so normalising before computing $E$ avoids getting $2E$.
By convention $\ket{1}$ is the $+E$ state, so the returned $\psi$ is always the positive-energy member; $C\psi$ is its
$-E$ partner.

\textbf{Step 5: gap check (optional).} If \code{gap\_factor} is set, \code{lowest\_bulk} gives $|E_{\text{bulk}}|$,
and a warning is issued if $|E_{\text{bulk}}| < \text{gap\_factor}\cdot E$, i.e.\ if the doublet does not sit well
below the bulk. Otherwise \code{E\_bulk = NaN}.

\textbf{Returns} \code{(ψ, E, γL, γR, E\_bulk)}. \textbf{Used} in cell 9 for the $\ell=6$ chain \code{H0} and inside
\code{charge\_difference} and \code{charge\_difference\_method\_2}. The time-dependent part uses a copy without the
energy and gap steps, \code{initial\_mode}.
\end{function}

\subsection{Static setup and profile helpers (cells 8--12)}

Cell 8 defines \code{Lℓ = 6}, \code{Nℓ = 64}, \code{V\_in, V\_out = -1.0, -4.0} and
\begin{Verbatim}
region(ℓ) = ((Nℓ - ℓ) ÷ 2 + 1):((Nℓ - ℓ) ÷ 2 + ℓ)
Vf(n) = (n + 1) in reg ? V_in : V_out
\end{Verbatim}
\code{region(ℓ)} is the centred range of $\ell$ 1-indexed sites. \code{Vf} receives the 0-indexed site from
\code{add\_onsite!}, hence \code{n + 1}. The notebook sets \code{ℓ = 6}, so \code{reg = 30:35}: in \code{H0} the six sites
$x=30\dots35$ are topological, with the Majoranas at the two ends of that region. Cell 9 calls \code{fermion\_mode(H0.mpo, H0.sites, Lℓ)}.

\begin{function}{site\_weights(γ), C(ψ), τx\_s}
\begin{Verbatim}
site_weights(γ) = [sum(abs2(inner(basis_state(H0.sites, vcat([nb], to_bits(x - 1, H0.L))), γ)) for nb in 0:1)
                   for x in 1:H0.N]
\end{Verbatim}
\code{site\_weights} returns $w(x) = |u(x)|^2+|v(x)|^2$ for $x=1,\dots,N$: the probability density of a BdG state
summed over the Nambu index. It gives the profiles of $\gamma_L$, $\gamma_R$, $\psi$ and $\bar\psi=C\psi$ plotted in
Figure~\ref{fig:majorana}. \code{τx\_s} and \code{C} in cell 10 are global copies of the particle--hole operator
defined inside \code{fermion\_mode}.
\end{function}

\begin{figure}[H]
\centering
\includegraphics[width=0.49\textwidth]{cell11_1.png}\hfill
\includegraphics[width=0.49\textwidth]{cell12_1.png}
\caption{Left: $|\gamma_L(x)|^2$ and $|\gamma_R(x)|^2$ from \code{site\_weights}. Right: $|\psi(x)|^2$ for
$(\gamma_L\pm i\gamma_R)/\sqrt2$, i.e.\ $\psi_f$ and $C\psi_f$, which have identical profiles. Both panels are for the
notebook's \code{H0} with $\ell=6$ (sites 30--35 topological): $\gamma_L$ peaks at $x=30$ and $\gamma_R$ at $x=35$, the
two ends of the topological region, and $\psi_f$ has weight at both.}
\label{fig:majorana}
\end{figure}

%======================================================================
\section{KPM local density of states and the gap estimate}\label{sec:kpm}
%======================================================================

\paragraph{LDOS call (cell 13).}
\begin{Verbatim}
Ncheb = 200; mdim = 60
omega = range(-4.0, 4.0; length=1500)
ldos = TensorBinding.get_ldos_spatial(H0, Ncheb, collect(omega); num_x=H0.N, mode=:mpo, kernel=:jackson,
                                      maxdim=mdim, verbose=false)
\end{Verbatim}
The library routine computes the spatially resolved LDOS $A(x,\omega)$ with the kernel polynomial method (KPM):
200 Chebyshev moments, MPO mode, bond dimension 60, on 1500 energies in $[-4,4]$ for all 64 sites. The printed log
shows that the spectral bounds were estimated by DMRG ($E_{\min}=-5.990$, $E_{\max}=5.990$, scale $6.589$); the
Hamiltonian is rescaled to $[-1,1]$ with these bounds. The Chebyshev moments are damped with the Jackson kernel
(the library default, written out in the call; Lorentz, Fej\'er and Dirichlet kernels are the alternatives); see
Weiße \emph{et al.}, Rev.\ Mod.\ Phys.\ \textbf{78}, 275 (2006). Because \code{H0} has a Nambu index, the library
switches the Nambu projection on automatically (printed info line); with the sector selector \code{proj\_nambu} left
at \code{nothing} it sums the particle and hole sectors, so
$A(x,\omega)=\sum_n\bigl(|u_n(x)|^2+|v_n(x)|^2\bigr)\,\delta_{\text{KPM}}(\omega-E_n)$, symmetric in $\omega$. The
result is a matrix with one row per energy and one column per site; it is plotted in Figure~\ref{fig:ldos}.

\begin{figure}[H]
\centering
\includegraphics[width=0.62\textwidth]{cell13_1.png}
\caption{KPM LDOS $A(x,\omega)$ of \code{H0} ($\ell=6$). The near-zero pair ($E_{\text{mode}}=0.019$) appears at
$x=30$ and $x=35$, the ends of the topological region; states of the region itself lie at $|\omega|\gtrsim1.2$, and the
bulk bands of the trivial chain start near $|\omega|\approx2$.}
\label{fig:ldos}
\end{figure}

\begin{function}{kpm\_gap(ldos, omega, E\_mode; mode\_width=0.3)}
\begin{Verbatim}
ω = collect(omega)
dos = vec(sum(ldos; dims=2))
D(w) = (dos[argmin(abs.(ω .- w))] + dos[argmin(abs.(ω .+ w))]) / 2
ws = ω[ω .> abs(E_mode) + mode_width]
Ds = D.(ws)
noise = 0.1 * maximum(Ds)
j = findfirst(k -> Ds[k] > noise && Ds[k] >= Ds[k+1], 1:length(Ds)-1)     # first bulk peak
half = Ds[j] / 2
i = findlast(k -> Ds[k] < half, 1:j)                                       # rising flank
return ws[i] + (half - Ds[i]) * (ws[i+1] - ws[i]) / (Ds[i+1] - Ds[i])
\end{Verbatim}
\textbf{Computes} the bulk gap edge $E_{\text{gap}}$ from the KPM LDOS.

\textbf{Method.}
\begin{enumerate}
\item Total DOS: $\mathrm{DOS}(\omega)=\sum_x A(x,\omega)$, symmetrised as
      $D(w)=[\mathrm{DOS}(w)+\mathrm{DOS}(-w)]/2$ (BdG symmetry), evaluated on the nearest grid points.
\item Only energies $w > |E_{\text{mode}}| + \text{mode\_width}$ are kept, so the mode's own peak is not mistaken for
      the band edge. The comment sets \code{mode\_width} $\approx 3\times$ the KPM resolution
      $\pi\cdot\text{scale}/N_{\text{cheb}}\approx 0.1$.
\item The first bulk peak is the first local maximum ($D_k\ge D_{k+1}$) above a noise floor of $10\%$ of the maximum.
\item $E_{\text{gap}}$ is the half-maximum point on the rising flank of that peak, found by linear interpolation
      between the last grid point below half maximum and the next one.
\end{enumerate}
The half-maximum is used because, per the comment, the KPM-broadened peak sits about 0.1 above the true edge.

\textbf{Returns} $E_{\text{gap}}$ (a \code{Float64}). \textbf{Used} in cell 16 to set the Fermi level. For \code{H0}
it gives $E_{\text{gap}}=1.154$; the DMRG value from \code{lowest\_bulk} for the same chain ($\ell=6$ in the sweep of
\S\ref{sec:dq}) is $1.275$.
\end{function}

%======================================================================
\section{Bulk projector and logical states}\label{sec:gamma}
%======================================================================

Cell 16 builds $P_{\text{bulk}}$, $\Gamma_0$ and $\Gamma_1$ as MPOs and checks them. It defines no named function
apart from two one-line helpers (\code{Cop}, \code{pbasis}); the steps are described in order.

\paragraph{Fermi level.} $P_{\text{bulk}}$ must contain all negative-energy bulk states and neither member of the
$\pm E$ mode. The Fermi level is therefore placed midway in the window between the bulk edge and the lower member of
the mode:
\[
  \epsilon_F = -\frac{E_{\text{mode}} + E_{\text{gap}}}{2},
\]
with $E_{\text{mode}}$ the energy returned by \code{fermion\_mode} and $E_{\text{gap}}$ from \code{kpm\_gap}. Printed:
$E_{\text{mode}}=0.018763$, $E_{\text{gap}}=1.15365$, $\epsilon_F=-0.58621$.

\paragraph{McWeeny purification.}
\begin{Verbatim}
ρ_bulk = TensorBinding.mcweeny_purify(H0; ϵF=ϵF, maxiters=60, maxdim=200, cutoff=1e-20, tol=1e-8)
\end{Verbatim}
McWeeny purification iterates $\rho\to3\rho^2-2\rho^3$, which sends eigenvalues above $\tfrac12$ to 1 and those below
to 0 (the comment in \code{mcweeny\_bulk}, \S\ref{sec:td}). Starting from a linear guess in $H$ whose eigenvalues
exceed $\tfrac12$ exactly for $E<\epsilon_F$, the result is the projector
$P_{\text{bulk}}=\sum_{\text{bulk }E_n<0}\ket{\psi_n}\bra{\psi_n}$. From the package source, the
\code{TBHamiltonian} method uses the initial guess
$\rho_0 = \bigl(\tfrac12 + \tfrac{\epsilon_F+\text{center}}{2\,\text{scale}}\bigr)\mathbb 1 - \tfrac{H}{2\,\text{scale}}$
with the Hamiltonian's stored KPM \code{center} and \code{scale}.

\textbf{Accuracy settings.} The markdown states that accuracy needs \code{cutoff=1e-20} and \code{tol=1e-8}. With
looser settings, or by iterating past convergence, $P_{\text{bulk}}$ drifts to an error of about $10^{-5}$, because
McWeeny does not correct errors that mix occupied and empty states, so they accumulate.

\paragraph{Logical states.}
\begin{Verbatim}
τx = ITensor([0.0 1.0; 1.0 0.0], H0.sites[1]', H0.sites[1])
Cop(ψ) = (φ = dag(ψ); φ[1] = noprime(τx * φ[1]); φ)
psi_bar = Cop(psi)
Γ0 = +(ρ_bulk, outer(psi_bar', psi_bar); cutoff=1e-20)
Γ1 = +(ρ_bulk, outer(psi', psi); cutoff=1e-20)
\end{Verbatim}
\code{Cop} is again $C=\tau_xK$. \code{outer(ψ', ψ)} is the rank-one MPO $\ket{\psi}\bra{\psi}$.

\paragraph{Checks (printed).}
\begin{itemize}
\item $\Tr P_{\text{bulk}} = 62.999999999993$, equal to the rank $N-1=63$ expected for the projector.
\item Mode occupation $\bra{\psi}\Gamma\ket{\psi}$: $9.0\times10^{-13}$ for $\ket{0}$, $1.0000000000009$ for $\ket{1}$.
      $\bra{\bar\psi}\Gamma\ket{\bar\psi}$: $0.99999999999997$ for $\ket{0}$, $-1.7\times10^{-14}$ for $\ket{1}$.
\item Many-body energies $E_{0,1}=\tfrac12\Tr(H\Gamma_{0,1})$ (up to a constant), evaluated exactly as
      \code{inner(H0.mpo, Γ)/2} since \code{inner(A,B)} $=\Tr(A^\dagger B)$:
      $E_1-E_0 = 0.0187632080507$, compared with $E=0.0187632080507$ from \code{fermion\_mode}. The comment notes
      that this difference is small only when the MZMs are far apart; here $\ell=6$.
\end{itemize}

\begin{function}{pbasis(x)}
\begin{Verbatim}
pbasis(x) = basis_state(H0.sites, vcat([0], to_bits(x - 1, H0.L)))
n0 = [real(inner(pbasis(x)', Γ0, pbasis(x))) for x in 1:H0.N]
\end{Verbatim}
The particle-sector basis state at site $x$ (the BdG vector of $c_x$). It is used to read the densities
$n_{0,1}(x)=\langle 0,x\vert\Gamma_{0,1}\vert 0,x\rangle$ shown in Figure~\ref{fig:dens}. A code comment notes that
the difference $\delta n$ should be computed directly from $\psi_f$, Eq.~\eqref{eq:dq}, rather than by subtracting two
purified MPOs.
\end{function}

\begin{figure}[H]
\centering
\includegraphics[width=0.6\textwidth]{cell16_1.png}
\caption{Local densities $n_0(x)$ and $n_1(x)$ of the two logical states of \code{H0} ($\ell=6$). They differ only
in the topological region $x=30\dots35$; the total difference is $\Delta Q=-0.105$ (the $\ell=6$ entry of
\S\ref{sec:dq}).}
\label{fig:dens}
\end{figure}

%======================================================================
\section{Charge difference versus length of the topological region}\label{sec:dq}
%======================================================================

One topological region ($V=-1$) of length $\ell$ is centred in the 64-site chain, with $V=-4$ elsewhere. The MZMs sit
at the two ends of the region, so shrinking $\ell$ pushes them together. For each $\ell$ the code comment lists the
pipeline: zero-mode pair from penalized DMRG on $K$ $\to$ gap check (third DMRG) $\to$ localise with $\hat X$ $\to$
self-conjugate $\gamma$ $\to$ pair into $\psi_f$ $\to$ keep the $+E$ member $\to$ $\Delta Q=\bra{\psi_f}\tau_z\ket{\psi_f}$.
Cell 18 repeats the definitions of \code{Lℓ}, \code{Nℓ}, \code{V\_in}, \code{V\_out} and \code{region}.

\begin{function}{τz\_on(ψ, s), ΔQ\_mode(ψ, s)}
\begin{Verbatim}
τz_on(ψ, s) = (φ = copy(ψ); φ[1] = noprime(ITensor([1.0 0.0; 0.0 -1.0], s[1]', s[1]) * φ[1]); φ)
ΔQ_mode(ψ, s) = real(inner(ψ, τz_on(ψ, s)))
\end{Verbatim}
\code{τz\_on} applies $\tau_z$ to the Nambu site. \code{ΔQ\_mode} returns
$\bra{\psi}\tau_z\ket{\psi}=\sum_x(|u(x)|^2-|v(x)|^2)$, which for $\psi=\psi_f$ is the charge difference
$\Delta Q$ of Eq.~\eqref{eq:dq}. The site list \code{s} only supplies the Nambu index. Since $C$ maps
$(u,v)\to(v^*,u^*)$, $\Delta Q(C\psi)=-\Delta Q(\psi)$. Used by \code{charge\_difference} and throughout the
time-dependent section.
\end{function}

\begin{function}{charge\_difference(ℓ; η=1e-2, seed=1234, cutoff=1e-20, gap\_factor=2.0)}
\begin{Verbatim}
reg = region(ℓ)
Vf(n) = (n + 1) in reg ? V_in : V_out
H = TensorBinding.get_Hamiltonian("chain_1d", 1.0; L = Lℓ)
TensorBinding.add_onsite!(H, Vf); TensorBinding.add_superconductivity!(H, 1.0)
pb(x) = basis_state(s, vcat([0], to_bits(x - 1, Lℓ)))
onsite_err = maximum(abs(real(inner(pb(x)', H.mpo, pb(x))) - Vf(x - 1)) for x in 1:Nℓ)
onsite_err > 1e-6 && @warn "ℓ = $ℓ: on-site potential of H deviates from Vf by $onsite_err"
psi, E, γL, γR, E_bulk = fermion_mode(H.mpo, s, Lℓ; η=η, seed=seed, cutoff=cutoff, gap_factor=gap_factor)
return (ΔQ = ΔQ_mode(psi, s), E = E, E_bulk = E_bulk, psi = psi, s = s)
\end{Verbatim}
\textbf{Method 1.} Builds the Hamiltonian for region length $\ell$, checks that the diagonal particle element of the
MPO reproduces $V(x)$ at every site to $10^{-6}$ (the MPO of the on-site term is compressed by the library), calls
\code{fermion\_mode} with the gap check enabled (\code{gap\_factor = 2}: warn if $|E_{\text{bulk}}|<2E$), and returns
a named tuple $(\Delta Q, E, E_{\text{bulk}}, \psi_f, s)$ with $\Delta Q$ for the $+E$ member ($\ket{1}$). The
mode \code{psi} and its site list \code{s} are kept so that \code{charge\_difference\_method\_2} can reuse them.

\textbf{Used} in cell 20 for $\ell = 1,\dots,40$; the cell prints the table $(\ell,\Delta Q,E,|E_{\text{bulk}}|)$
of \S\ref{sec:results}. Apart from that table the output contains only the library's $p$-wave info lines, no on-site
or gap-check warnings.
\end{function}

\begin{function}{charge\_difference\_method\_2(r)}
\begin{Verbatim}
psi, s = r.psi, r.s
psibar = C(psi)
upart(ψ) = [abs2(inner(pb(x),ψ)) for x in 1:Nℓ]
n_density = upart(psi) - upart(psibar)
charge_diff = sum(n_density)
\end{Verbatim}
\textbf{Method 2.} Takes the entry \code{r} of the $\ell$ sweep (\code{sweep\_ℓ}), i.e.\ the same DMRG mode $\psi_f$
as method 1, so no second DMRG is run, and evaluates
$\Delta Q$ as the sum of the density difference read off the projectors,
$\delta n(x)=|u_f(x)|^2-|u_{\bar f}(x)|^2$ with $\bar\psi_f=C\psi_f$ (the comment calls this ``using projectors
$\Gamma_0$ and $\Gamma_1$''; the bulk cancels, so only $\psi_f$ and $C\psi_f$ enter). By Eq.~\eqref{eq:dq} it equals
method 1 analytically: $u_{\bar f}=v_f^*$, so $\sum_x(|u_f|^2-|u_{\bar f}|^2)=\bra{\psi_f}\tau_z\ket{\psi_f}$. Method 2
is therefore an independent evaluation of the same number. It returns a single number. Used in cell 24 for every
entry of the sweep; the cell prints $\max_\ell|\Delta Q_\ell^{(1)}-\Delta Q_\ell^{(2)}| = 5.6\times10^{-16}$, and the
two curves coincide in Figure~\ref{fig:dql}.
\end{function}

\begin{figure}[H]
\centering
\includegraphics[width=0.62\textwidth]{cell25_1.png}
\caption{$|\Delta Q|$ (units of $-e$) versus the topological region length $\ell$, axis reversed so that the MZMs
approach each other from left to right. Method 1 (\code{charge\_difference}) and method 2
(\code{charge\_difference\_method\_2}, dashed) are plotted on top of each other and agree to $6\times10^{-16}$.}
\label{fig:dql}
\end{figure}

%======================================================================
\section{Time-dependent ramp}\label{sec:td}
%======================================================================

\subsection{Protocol}
From the markdown: at $t=0$ the 32-site chain ($L_t=5$) is topological ($\mu_{\text{topo}}=-0.4$) with MZMs at its
ends. Sites are ramped to $\mu_{\text{triv}}=-5.6$ from the chain ends inward, each with a small delay, which carries
the MZMs to the middle. At $t=T$ only the centre site ($\ell_{\text{center}}=1$, $x=16$) is topological, so the two
MZMs overlap and the logical states differ in charge.
\[
  \mu(x,t) = \mu_{\text{topo}} + (\mu_{\text{triv}}-\mu_{\text{topo}})\,s\bigl(q_i(t)\bigr),\qquad
  q_i(t) = \frac tT\bigl[1+\alpha(N_{\text{leg}}-1)\bigr]-\alpha i,\qquad
  s(\tau)=3\tau^2-2\tau^3 ,
\]
with $i$ the distance of $x$ from the nearest chain end, $\alpha=0.025$ and $s$ clamped to $[0,1]$. The
$N_t-\ell_{\text{center}}=31$ ramped sites cannot split evenly: the left leg has 15 sites ($x=1\dots15$), the right leg 16
($x=17\dots32$), and $N_{\text{leg}}=16$ is the longer leg.
\begin{itemize}
\item Site-delayed ramp: site $i$ starts at $t=\alpha iT/[1+\alpha(N_{\text{leg}}-1)]$ and ramps over the duration
      $T/[1+\alpha(N_{\text{leg}}-1)]$; the trivial region grows inward from both ends at the same rate. The innermost site
      of the right leg ($i=15$) finishes at $t=T$, that of the left leg ($i=14$) at $t/T=(1+14\alpha)/(1+15\alpha)\approx0.98$.
\item A uniform ramp would cross $|\mu|=2t$ everywhere at once and close the bulk gap. With the delay the instantaneous
      bulk gap stays $\ge0.83$ (minimum near $t/T\approx0.43$).
\item The smoothstep starts and stops with zero velocity.
\end{itemize}
\textbf{What is evolved.} $\Gamma_1(t)-\Gamma_0(t) = \ket{\psi_f}\bra{\psi_f}-\ket{C\psi_f}\bra{C\psi_f}$ at all times,
so $\Delta Q(t)=\bra{\psi_f(t)}\tau_z\ket{\psi_f(t)}$ needs only the single BdG vector $\psi_f$, evolved with
$i\partial_t\psi=H(t)\psi$ by TDVP. At $t=T$ the weight of $\psi_f$ outside the instantaneous $\pm E$ doublet is the
leakage into bulk quasiparticles (the non-adiabatic error). Full densities add the shared bulk, either TDVP-evolved
(all $N_t-1$ filled states) or the adiabatic McWeeny projector of $H(t)$.

\subsection{Initial mode and final branches (cell 27)}

\begin{function}{C\_t(ψ)}
\begin{Verbatim}
C_t(ψ) = (φ = dag(ψ); s = siteind(φ, 1); φ[1] = noprime(ITensor([0.0 1.0; 1.0 0.0], s', s) * φ[1]); φ)
\end{Verbatim}
Particle--hole conjugation $C=\tau_xK$ that reads the Nambu index from the MPS's own first site, so it does not
depend on a global site list. The comment states that the functions of this cell work for any chain length and read
no globals.
\end{function}

\begin{function}{initial\_mode(Hmpo, sites, L)}
Same as steps 1--4 of \code{fermion\_mode}: \code{doublet}, localisation with $\hat X$, \code{fix} for
$C\gamma=\gamma$ (using \code{C\_t}), and $\psi=(\gamma_L+i\gamma_R)/\lVert\cdot\rVert$. Differences: it uses the
ordinary eigenproblem \code{eigen(Hermitian((X + X') / 2))} (the doublet is already orthonormal, so $S=\mathbb 1$ up
to round-off), it does not compute $E$ and does not choose the $+E$ branch, and it uses the default
$\eta$, seed and cutoff of \code{doublet}. Returns the normalised $\psi$.

Used once (cell 29): \code{ψf\_init = initial\_mode(H\_of\_t\_alt(0.0, 1.0, smoothstep), st, Lt)}; $H(0)$ is the same
for every $T$ and ramp. Printed checks: $\lVert H\psi\rVert = 1.12\times10^{-12}$ (the MZMs at the ends of the
32-site topological chain are at zero energy to this precision) and $\Delta Q(0)=4.6\times10^{-15}$.
\end{function}

\begin{function}{final\_branches(Hmpo, sites)}
\begin{Verbatim}
ψ1, ψ2 = doublet(Hmpo, sites)
Hm = [inner(a', Hmpo, b) for a in ψs, b in ψs]
F = eigen(Hermitian((Hm + Hm') / 2))
mk(c) = add2(c[1] * ψ1, c[2] * ψ2)
return mk(F.vectors[:, 2]), mk(F.vectors[:, 1]), F.values[2]    # ψ₊, ψ₋, E₊
\end{Verbatim}
Instantaneous $\pm E$ eigenstates at the end of the ramp. The $2\times2$ matrix of $H$ in the doublet is
diagonalised; the comment notes that $E$ is large at $t=T$, so $H$ itself separates the two states (no localisation
is needed). Returns $(\psi_+,\psi_-,E_+)$. Used by \code{run\_ramp}, \code{evolve\_states} and
\code{evolve\_states\_single\_modes} to project the evolved state and to identify which logical state ended in $+E$.
\end{function}

\subsection{Time-dependent Hamiltonian (cell 28)}

\begin{function}{smoothstep(τ)}
\begin{Verbatim}
smoothstep(τ) = τ <= 0 ? 0.0 : τ >= 1 ? 1.0 : τ^2 * (3 - 2τ)
\end{Verbatim}
$s(\tau)=3\tau^2-2\tau^3$ clamped to $[0,1]$; $s'(0)=s'(1)=0$.
\end{function}

\begin{function}{mu\_xt(x, t, T, ramp; α = alpha)}
\begin{Verbatim}
ℓ_center = 1                              # sites kept topological in the middle
N_left  = (Nt - ℓ_center) ÷ 2             # 15 ramped sites on the left  (x = 1 … 15)
N_right = Nt - ℓ_center - N_left          # 16 ramped sites on the right (x = 17 … 32)
N_leg   = max(N_left, N_right)            # 16
function mu_xt(x, t, T, ramp; α = alpha)
    (x <= N_left || x > Nt - N_right) || return μ_topo_t
    i = min(x - 1, Nt - x)
    q = (t/T) * (1 + α * (N_leg - 1)) - α * i
    return μ_topo_t + (μ_triv_t - μ_topo_t) * ramp(q)
end
\end{Verbatim}
Implements $\mu(x,t)$ of the protocol ($\ell_{\text{center}}=1$, \code{Nt = 32}, $\alpha=0.025$). A site is ramped if it
lies in the left leg ($x\le$ \code{N\_left}) or the right leg ($x>N_t-$ \code{N\_right}); its delay index is the
distance from the nearer chain end, $i=\min(x-1,N_t-x)$, so the two walls move inward at the same rate. With
$\ell_{\text{center}}=1$ the 31 ramped sites split into 15 (left, $i=0..14$) and 16 (right, $i=0..15$), and
$N_{\text{leg}}=\max(15,16)=16$ makes the innermost site of the longer leg finish exactly at $t=T$; site 16 is never
ramped. The same function works for even $\ell_{\text{center}}$ (equal legs). The argument \code{ramp} is the profile
function (\code{smoothstep}).
\end{function}

\begin{function}{onsite\_mpo(V)}
\begin{Verbatim}
function onsite_mpo(V)
    arr = zeros(ntuple(_ -> 2, Lt)...)
    for I in CartesianIndices(arr)
        b = Tuple(I) .- 1   # position bits, most significant first
        arr[I] = V(1 + sum(b[k] * 2^(Lt - k) for k in 1:Lt))
    end
    D = TensorBinding.mps_to_diagonal_mpo(MPS(arr, pos_t; cutoff=1e-16), pos_t)
    return TensorBinding.prepend_nambu(D, nambu_t, :tz)
end
\end{Verbatim}
Builds the exact on-site MPO $\tau_z\otimes\operatorname{diag}V(x)$ from the $N_t$ values $V(x)$. The values are
written into a dense $2\times\dots\times2$ array indexed by the position bits, compressed to an MPS on the position
sites (\code{pos\_t = st[2:end]}) with cutoff $10^{-16}$, turned into a diagonal MPO, and given the Nambu structure
$\tau_z$ on \code{nambu\_t = st[1]}. Per the comment, this avoids rebuilding the whole Hamiltonian (and running QTCI)
at every time step.
\end{function}

\begin{function}{H\_of\_t\_alt(t, T, ramp), pbt(x)}
\begin{Verbatim}
H_of_t_alt(t, T, ramp) = +(H_static.mpo, onsite_mpo(x -> mu_xt(x, t, T, ramp)); cutoff=1e-16)
pbt(x) = basis_state(st, vcat([0], to_bits(x - 1, Lt)))
\end{Verbatim}
$H(t) = H_{\text{static}} + \tau_z\otimes\operatorname{diag}\mu(x,t)$, where \code{H\_static} (hopping and pairing,
no on-site term) is built once with \code{get\_Hamiltonian("chain\_1d", 1.0; L = Lt)} and
\code{add\_superconductivity!(H\_static, 1.0)}; \code{st = H\_static.sites}. \code{pbt(x)} is the particle-sector
basis state at $x$. Cell 28 checks $\max_x|\langle0,x|H|0,x\rangle-\mu(x,t)|$ for $T=80$:
\begin{center}
\begin{tabular}{@{}rcc@{}}
\toprule
$t$ & max on-site error & \code{maxlinkdim(H)} \\
\midrule
0  & $1.8\times10^{-15}$ & 5 \\
40 & $3.4\times10^{-14}$ & 8 \\
80 & $1.1\times10^{-14}$ & 6 \\
\bottomrule
\end{tabular}
\end{center}
\end{function}

\begin{function}{norm\_Hψ(H, ψ)}
\begin{Verbatim}
norm_Hψ(H, ψ) = norm(orthogonalize(apply(H, ψ; alg="naive", truncate=false), 1)[1])
\end{Verbatim}
$\lVert H\psi\rVert$ computed without truncation: $H\psi$ is formed exactly, brought to orthogonality centre 1, and the
norm of the first tensor is the norm of the MPS. Used once to check $\psi_f(0)$.
\end{function}

\subsection{Evolution of the fermion mode}

\begin{function}{run\_ramp(T, ramp; dt=0.1, nsnap=100)}
\begin{Verbatim}
nsteps = round(Int, T / dt); dt = T / nsteps
states = TensorBinding.evolve_with_tdvp_timedep(t -> H_of_t_alt(t, T, ramp), ψf_init, nsteps, dt;
                                                maxdim=64, cutoff=1e-20, tol=1e-12, krylovdim=30,
                                                reverse_step=true)
ψ₊, ψ₋, E₊ = final_branches(H_of_t_alt(T, T, ramp), st)
ψT = states[end]
p₊, p₋ = abs2(inner(ψ₊, ψT)), abs2(inner(ψ₋, ψT))
σ = p₋ > p₊ ? -1 : 1
σ == -1 && ((p₊, p₋) = (p₋, p₊))
idx = unique(round.(Int, range(1, nsteps + 1; length=nsnap + 1)))
return (T = T, leak = 1 - p₊ - p₋, p_wrong = p₋, E₊ = E₊,
        ΔQ_final = σ * ΔQ_mode(ψT, st), ΔQ_adiabatic = ΔQ_mode(ψ₊, st),
        times = (idx .- 1) .* dt, ΔQ_t = [σ * ΔQ_mode(states[i], st) for i in idx], snapshots = states[idx])
\end{Verbatim}
\textbf{Computes} the evolution $i\partial_t\ket{\psi(t)}=H(t)\ket{\psi(t)}$ of $\psi_f(0)$ from $0$ to $T$ and the
resulting non-adiabatic error and $\Delta Q(t)$.

\textbf{Time stepping.} \code{dt} is adjusted so that \code{nsteps*dt = T}. From the package source, the routine
\code{evolve\_with\_tdvp\_timedep}\linebreak[3] evaluates $H$ at the step midpoint $t_{\text{mid}}=(k-\tfrac12)\,dt$, uses two-site
TDVP by default (\code{nsite = 2}), normalises after each step by default, and returns all \code{nsteps+1} states.
The notebook sets \code{maxdim = 64}, \code{cutoff = 1e-20}, Krylov tolerance $10^{-12}$ with Krylov dimension 30,
and \code{reverse\_step = true} (the comment: a sweep consists of a forward and a backward pass through the MPS).

\textbf{Analysis.}
\begin{itemize}
\item $p_\pm = |\langle\psi_\pm\vert\psi(T)\rangle|^2$ with the instantaneous branches from \code{final\_branches}.
\item Branch identification: if the state ended mainly in $\psi_-$, then $\sigma=-1$ and $p_\pm$ are swapped. Since
      $\Delta Q(C\psi)=-\Delta Q(\psi)$, multiplying by $\sigma$ reports $\Delta Q$ for the logical state that ends in
      $+E$ ($\ket{1}$).
\item \code{leak} $=1-p_+-p_-$: the probability transferred into bulk quasiparticles.
      \code{p\_wrong} $=p_-$ after the swap: the weight in the other branch.
\item \code{ΔQ\_final} $=\sigma\Delta Q(\psi(T))$; \code{ΔQ\_adiabatic} $=\Delta Q(\psi_+)$, the adiabatic value.
\item \code{times}, \code{ΔQ\_t}, \code{snapshots}: \code{nsnap+1} (after \code{unique}) roughly equally spaced
      samples.
\end{itemize}
\textbf{Used} in cell 31 for $T=150,200,250,300,350$, which prints \code{leak}, \code{p\_wrong}, \code{ΔQ\_final},
\code{ΔQ\_adiabatic} and their difference (results in \S\ref{sec:results}), and in cell 32 (Figure~\ref{fig:dqt}).
\end{function}

\begin{figure}[H]
\centering
\includegraphics[width=0.6\textwidth]{cell32_1.png}
\caption{$\Delta Q(t)$ of the evolved $\psi_f$ for the smoothstep ramp with $T=350$ (\code{run\_ramp}). $\Delta Q$
stays at zero while the MZMs are far apart and goes to $-0.889$ as they fuse at the centre. Dashed: the adiabatic
value \code{ΔQ\_adiabatic} $=\Delta Q(\psi_+)=-0.889108$ of the final instantaneous mode.}
\label{fig:dqt}
\end{figure}

\begin{function}{evolve\_states(T, ramp; dt=0.1, nsnap=160)}
\textbf{Computes} $\delta n(x,t)=n_1(x,t)-n_0(x,t)=|u_f(x,t)|^2-|u_{\bar f}(x,t)|^2$ for both logical states.
\begin{Verbatim}
s1, s0 = ev(ψf_init), ev(C_t(ψf_init))
ψ₊, _, _ = final_branches(H_of_t_alt(T, T, ramp), st)
abs2(inner(ψ₊, s0[end])) > abs2(inner(ψ₊, s1[end])) && ((s1, s0) = (s0, s1))   # |1⟩: mode ends in +E
upart(ψ) = [abs2(inner(pbt(x), ψ)) for x in 1:Nt]
δn = reduce(hcat, [upart(s1[i]) .- upart(s0[i]) for i in idx])
return (T = T, times = (idx .- 1) .* dt, δn = δn,
        C_err = norm(add2(C_t(s1[end]), -1 * s0[end])),
        n_mode = (abs2(inner(ψ₊, s1[end])), abs2(inner(ψ₊, s0[end]))))
\end{Verbatim}
\textbf{Method.} Both $\psi_f$ and $\bar\psi_f=C\psi_f$ are evolved by TDVP with the settings of \code{run\_ramp}. The
one whose final state has more weight on $\psi_+$ is labelled $\ket{1}$. Per the comments: the bulk cancels exactly in
$\Gamma_1-\Gamma_0$, and evolving and subtracting the two full $\Gamma$ MPOs would bury $\delta n$ (about $10^{-14}$
early on) under truncation error; $\bar\psi_f$ is evolved independently rather than taken as $C\psi_f(t)$, so that
$\lVert C\psi_f(T)-\bar\psi_f(T)\rVert$ (\code{C\_err}) checks the evolution. \code{n\_mode} holds the weights of
the two final states on $\psi_+$.

\textbf{Returns} \code{(T, times, δn, C\_err, n\_mode)}, with \code{δn} an $N_t\times n_{\text{snap}}$ matrix.
\textbf{Used} in cell 33 for $T=300$ and $T=100$ (Figure~\ref{fig:dnxt}); the cell prints \code{C\_err} and
\code{n\_mode} for both runs (\S\ref{sec:results}).
\end{function}

\begin{figure}[H]
\centering
\includegraphics[width=0.85\textwidth]{cell34_1.png}
\caption{Top: $\delta n(x,t)=n_1-n_0$ for $T=300$ and $T=100$ (\code{evolve\_states}). The density difference
appears at the centre site $x=16$ once the MZMs come together. Bottom: $\Delta Q(t)=\sum_x\delta n(x,t)$ versus $t/T$;
the two ramp durations give almost the same curve, ending near $-0.89$.}
\label{fig:dnxt}
\end{figure}

\subsection{Bulk states and full densities (cells 35--39)}

\begin{function}{bdg\_basis, dense\_bdg(Hmpo), hbt(x)}
\begin{Verbatim}
bdg_basis = [basis_state(st, vcat([nb], to_bits(x - 1, Lt))) for nb in 0:1 for x in 1:Nt]   # index nb·Nt + x
dense_bdg(Hmpo) = [inner(a', Hmpo, b) for a in bdg_basis, b in bdg_basis]
hbt(x) = basis_state(st, vcat([1], to_bits(x - 1, Lt)))
\end{Verbatim}
\code{bdg\_basis} lists the $2N_t=64$ BdG basis states, ordered as $(u_1,\dots,u_{N_t},v_1,\dots,v_{N_t})$, i.e.\
index $n_b N_t + x$. \code{dense\_bdg} converts an MPO to the dense $2N_t\times2N_t$ BdG matrix by matrix elements.
\code{hbt(x)} is the hole-sector basis state, used to read $|v(x)|^2$.

\textbf{Convention check.} The cell also builds \code{Hs\_d = dense\_bdg(H\_static.mpo)} and compares it with the
reference matrix $[h\ \Delta;\ \Delta^\dagger\ {-h^T}]$ of \S\ref{sec:model} (\code{h\_ref}, \code{Δ\_ref}); it prints
$h_{12}=1$, $\Delta_{12}=1$, $\Delta_{21}=-1$, $\Delta_{11}=-1.9\times10^{-17}$ and a maximum deviation of $3.6\times10^{-15}$.
\end{function}

\begin{function}{vec\_to\_mps(c; cutoff=1e-16)}
Inverse of the above: writes the dense BdG vector $c$ (ordering $n_bN_t+x$) into a $2^{L_t+1}$ complex array indexed
by [Nambu bit, position bits] and compresses it into an MPS on \code{st} with cutoff $10^{-16}$.
\end{function}

\begin{function}{initial\_bulk\_states(Hmpo)}
\begin{Verbatim}
Hd = dense_bdg(Hmpo)
F = eigen(Hermitian((Hd + Hd') / 2))
doublet_idx = sortperm(abs.(F.values))[1:2]
bulk_idx = [n for n in eachindex(F.values) if F.values[n] < 0 && n ∉ doublet_idx]
@assert length(bulk_idx) == Nt - 1
return [vec_to_mps(F.vectors[:, n]) for n in bulk_idx]
\end{Verbatim}
The filled bulk states at $t=0$: all negative-energy eigenstates of $H(0)$ except the near-zero doublet (the two
smallest $|E|$), $N_t-1=31$ of them. For $N_t=32$ the BdG matrix is only $64\times64$, so the comment notes that exact
diagonalisation is used and the eigenvectors are converted to MPS. Stored as \code{bulk\_init}. Printed check:
$\lVert P_{\text{bulk}}\psi_f\rVert^2 = \sum_b|\langle b\vert\psi_f\rangle|^2 = 1.8\times10^{-25}$ at $t=0$.
\end{function}

\begin{function}{bulk\_density\_tdvp(T, ramp; dt=0.1, nsnap=160)}
\begin{Verbatim}
Hs = [H_of_t_alt((k - 0.5) * dt, T, ramp) for k in 1:nsteps]   # H at t_mid of step k
Hoft(t) = Hs[round(Int, t / dt + 0.5)]
...
s = TensorBinding.evolve_with_tdvp_timedep(Hoft, b, nsteps, dt; maxdim=64, cutoff=1e-20, tol=1e-12,
                                           krylovdim=30, reverse_step=true)
n_bulk[:, j] .+= upart(s[i]);  v_bulk[:, j] .+= vpart(s[i])
...
G = [inner(a, b) for a in final, b in final]
return (n_bulk = n_bulk, v_bulk = v_bulk, orth_err = maximum(abs, G - I))
\end{Verbatim}
\textbf{Computes} the evolved bulk density $n_{\text{bulk}}(x,t)=\sum_b|u_b(x,t)|^2$ by evolving each of the 31 bulk
states with the same $H(t)$ and TDVP settings as $\psi_f$. $H(t_{\text{mid}})$ for every step is built once and shared
by all states; \code{Hoft} maps the midpoint time $(k-\frac12)dt$ requested by the library back to index $k$. It also
accumulates the hole weight $v_{\text{bulk}}(x,t)=\sum_b|v_b(x,t)|^2$ (for the completeness check) and returns the
overlap error $\max_{mn}|\langle b_m\vert b_n\rangle-\delta_{mn}|$ of the final states, which should stay near zero.
Samples are at the same snapshot indices as \code{evolve\_states\_single\_modes}.
\end{function}

\begin{function}{evolve\_states\_single\_modes(T, ramp; dt=0.1, nsnap=160)}
Same evolution and $\ket{1}$ labelling as \code{evolve\_states}, but returns the two particle densities separately,
\code{n\_pos} $=|u_f(x,t)|^2$ (state $\ket{1}$) and \code{n\_neg} $=|u_{\bar f}(x,t)|^2$ (state $\ket{0}$), together
with \code{times}, \code{C\_err} and \code{n\_mode}. Used by \code{evolve\_full\_densities}.
\end{function}

\begin{function}{ramp\_fermi\_level(ramp; npts=401)}
\begin{Verbatim}
H_static_dense = dense_bdg(H_static.mpo)
E = map(range(0, 1; length=npts)) do τ
    μ = [mu_xt(x, τ, 1.0, ramp) for x in 1:Nt]
    sort(abs.(eigvals(Hermitian(H_static_dense + Diagonal(vcat(μ, -μ))))))   # E[1] = E[2] = mode, E[3] = bulk edge
end
Em, Eg = maximum(e[1] for e in E), minimum(e[3] for e in E)
Em < Eg || error("no single ϵF: max E_mode = $Em ≥ min E_gap = $Eg")
return -(Em + Eg) / 2, Em, Eg, maximum(e[end] for e in E)
\end{Verbatim}
\textbf{Computes} one Fermi level valid for the whole ramp. Because $\mu(x,t)$ depends only on $t/T$, the
instantaneous spectrum is scanned once over $t/T\in[0,1]$ (401 points) by exact diagonalisation of
$H(t)=H_{\text{static}}+\tau_z\otimes\operatorname{diag}\mu(x,t)$, here \code{Diagonal(vcat(μ, -μ))} in the
\code{bdg\_basis} ordering. $\epsilon_F$ must lie below $-E_{\text{mode}}(t)$ and above $-E_{\text{gap}}(t)$ at every
$t$, so the worst cases are used:
\[
  \epsilon_F = -\frac{\max_t E_{\text{mode}}(t) + \min_t E_{\text{gap}}(t)}{2}.
\]
Returns $(\epsilon_F, \max_t E_{\text{mode}}, \min_t E_{\text{gap}}, \max|E|)$. Printed: $\max_tE_{\text{mode}}=0.34875$,
$\min_tE_{\text{gap}}=0.82787$, window $(-0.82787,-0.34875)$, $\epsilon_F=-0.58831$, $\max|E|=7.5749$.
\end{function}

\begin{function}{mcweeny\_bulk(Hmpo; ϵF=ϵF\_ramp, scale=9.0, maxiters=80, maxdim=200, cutoff=1e-20, tol=1e-10)}
\begin{Verbatim}
Id = MPO(st, "Id")
ρ0 = +((0.5 + ϵF / (2scale)) * Id, (-0.5 / scale) * Hmpo; cutoff=cutoff)
ρ = TensorBinding.mcweeny_purify(ρ0; maxiters=maxiters, maxdim=maxdim, cutoff=cutoff, tol=tol)
abs(real(tr(ρ)) - (Nt - 1)) > 1e-6 && @warn "Tr P_bulk = $(real(tr(ρ))), expected $(Nt - 1): check ϵF"
return [real(inner(pbt(x)', ρ, pbt(x))) for x in 1:Nt]
\end{Verbatim}
\textbf{Computes} the instantaneous (adiabatic) bulk density $\langle0,x\vert P_{\text{bulk}}(t)\vert0,x\rangle$ for a
given $H(t)$. The initial guess
\[
  \rho_0 = \tfrac12 - \frac{H-\epsilon_F}{2\,\text{scale}}
\]
has eigenvalues in $[0,1]$ when $\text{scale}\ge\max|E-\epsilon_F|$; the comment cites $\max|E|\approx7.6$, and
scale $=9$ is used. McWeeny $\rho\to3\rho^2-2\rho^3$ drives eigenvalues above $\frac12$ (i.e.\ $E<\epsilon_F$) to 1 and
the rest to 0, which gives $P_{\text{bulk}}(t)$ when $\epsilon_F$ lies between the bulk edge and the lower member of
the doublet. A warning is issued if $\Tr\rho$ differs from $N_t-1$ by more than $10^{-6}$.
\end{function}

\begin{function}{bulk\_density\_mcweeny(T, ramp, times; kw...)}
\begin{Verbatim}
bulk_density_mcweeny(T, ramp, times; kw...) = reduce(hcat, [mcweeny_bulk(H_of_t_alt(t, T, ramp); kw...) for t in times])
\end{Verbatim}
Applies \code{mcweeny\_bulk} at each snapshot time and stacks the results into an $N_t\times n_{\text{snap}}$ matrix.
\end{function}

\begin{function}{evolve\_full\_densities(T, ramp; bulk=:tdvp, dt=0.1, nsnap=160)}
\begin{Verbatim}
m = evolve_states_single_modes(T, ramp; dt=dt, nsnap=nsnap)
if bulk == :tdvp
    b = bulk_density_tdvp(T, ramp; dt=dt, nsnap=nsnap)
    compl_err = maximum(abs, b.n_bulk .+ b.v_bulk .+ m.n_pos .+ m.n_neg .- 1)
elseif bulk == :mcweeny
    n_bulk, orth_err, compl_err = bulk_density_mcweeny(T, ramp, m.times), NaN, NaN
...
return (T = T, times = m.times, n1 = n_bulk .+ m.n_pos, n0 = n_bulk .+ m.n_neg, n_bulk = n_bulk,
        orth_err = orth_err, compl_err = compl_err, C_err = m.C_err, n_mode = m.n_mode)
\end{Verbatim}
\textbf{Computes} the full densities $n_1(x,t)=n_{\text{bulk}}+|u_f|^2$ and $n_0(x,t)=n_{\text{bulk}}+|u_{\bar f}|^2$,
with the shared bulk either TDVP-evolved (\code{:tdvp}) or adiabatic (\code{:mcweeny}). For \code{:tdvp} the
completeness error checks $\Gamma_1+C\Gamma_1C=\mathbb 1$, i.e.\ at every site
\[
  \sum_{\text{bulk}}\bigl(|u_b(x)|^2+|v_b(x)|^2\bigr)+|u_f(x)|^2+|v_f(x)|^2=1,
\]
using $|v_f|^2=|u_{\bar f}|^2$. Used in cell 39 with $T=100$ and \code{nsnap = 40} (Figure~\ref{fig:full}); the cell
prints \code{orth\_err}, \code{compl\_err}, \code{C\_err}, \code{n\_mode} and the TDVP--McWeeny bulk difference.
\end{function}

\begin{figure}[H]
\centering
\includegraphics[width=0.92\textwidth]{cell40_1.png}
\caption{Full densities for the smoothstep ramp with $T=100$ (\code{evolve\_full\_densities}, TDVP bulk). Top:
$n_1(x,t)$ and $n_0(x,t)$. Bottom left: $n_{\text{bulk}}(x,t)$ from TDVP. The trivial region (density near 1) grows
inward from both ends. Bottom right: $\max_x|n_{\text{bulk}}^{\text{TDVP}}-n_{\text{bulk}}^{\text{McWeeny}}|$ versus
$t$ (log scale), the difference between the evolved and the adiabatic bulk.}
\label{fig:full}
\end{figure}

%======================================================================
\section{Results}\label{sec:results}
%======================================================================

\subsection{\texorpdfstring{Static chain, $\ell=6$ (\code{H0})}{Static chain, l=6 (H0)}}
\begin{center}
\begin{tabular}{@{}ll@{}}
\toprule
Quantity & Value \\
\midrule
$E_{\text{mode}}$ (\code{fermion\_mode}) & 0.0187632081 \\
$E_{\text{gap}}$ (\code{kpm\_gap}) & 1.1536543281 \\
$\epsilon_F=-(E_{\text{mode}}+E_{\text{gap}})/2$ & $-0.5862087681$ \\
$\Tr P_{\text{bulk}}$ & 62.999999999993 \\
$\bra{\psi}\Gamma_0\ket{\psi}$, $\bra{\psi}\Gamma_1\ket{\psi}$ & $9.0\times10^{-13}$, 1.0000000000009 \\
$\bra{\bar\psi}\Gamma_0\ket{\bar\psi}$, $\bra{\bar\psi}\Gamma_1\ket{\bar\psi}$ & 0.99999999999997, $-1.7\times10^{-14}$ \\
$E_1-E_0=\tfrac12\Tr[H(\Gamma_1-\Gamma_0)]$ & 0.0187632081 \\
\bottomrule
\end{tabular}
\end{center}

\subsection{\texorpdfstring{$\Delta Q$ versus $\ell$ (\code{charge\_difference}, cell 20)}{Delta Q versus l (charge\_difference, cell 20)}}
The cell prints all 40 entries. $\Delta Q$ is negative for $\ket{1}$ (the
$+E$ member). $E$ halves with each added site from $\ell\approx5$ on, and $|\Delta Q|$ falls by about the same factor;
$|E_{\text{bulk}}|$ decreases from 2.02 to
1.006 and stays above $2E$ in every case (no gap-check warning).
\begin{center}\small
\begin{tabular}{@{}rlll@{\qquad}rlll@{}}
\toprule
$\ell$ & $\Delta Q$ & $E$ & $|E_{\text{bulk}}|$ & $\ell$ & $\Delta Q$ & $E$ & $|E_{\text{bulk}}|$\\
\midrule
1 & $-0.8481$ & $0.7379$ & 2.0181 & 21 & $-1.179\times10^{-5}$ & $5.722\times10^{-7}$ & 1.0235 \\
2 & $-0.6107$ & $0.322$ & 2.0181 & 22 & $-6.180\times10^{-6}$ & $2.861\times10^{-7}$ & 1.0214 \\
3 & $-0.4209$ & $0.1536$ & 1.7891 & 23 & $-3.233\times10^{-6}$ & $1.431\times10^{-7}$ & 1.0196 \\
4 & $-0.2766$ & $0.07557$ & 1.5330 & 24 & $-1.688\times10^{-6}$ & $7.153\times10^{-8}$ & 1.0179 \\
5 & $-0.1738$ & $0.03759$ & 1.3753 & 25 & $-8.798\times10^{-7}$ & $3.576\times10^{-8}$ & 1.0165 \\
6 & $-0.1052$ & $0.01876$ & 1.2746 & 26 & $-4.578\times10^{-7}$ & $1.788\times10^{-8}$ & 1.0152 \\
7 & $-0.06191$ & $9.377\times10^{-3}$ & 1.2077 & 27 & $-2.378\times10^{-7}$ & $8.941\times10^{-9}$ & 1.0141 \\
8 & $-0.03563$ & $4.688\times10^{-3}$ & 1.1616 & 28 & $-1.234\times10^{-7}$ & $4.470\times10^{-9}$ & 1.0131 \\
9 & $-0.02016$ & $2.344\times10^{-3}$ & 1.1288 & 29 & $-6.399\times10^{-8}$ & $2.235\times10^{-9}$ & 1.0122 \\
10 & $-0.01125$ & $1.172\times10^{-3}$ & 1.1048 & 30 & $-3.307\times10^{-8}$ & $1.118\times10^{-9}$ & 1.0114 \\
11 & $-6.211\times10^{-3}$ & $5.859\times10^{-4}$ & 1.0868 & 31 & $-1.710\times10^{-8}$ & $5.588\times10^{-10}$ & 1.0107 \\
12 & $-3.398\times10^{-3}$ & $2.930\times10^{-4}$ & 1.0730 & 32 & $-8.829\times10^{-9}$ & $2.794\times10^{-10}$ & 1.0100 \\
13 & $-1.846\times10^{-3}$ & $1.465\times10^{-4}$ & 1.0622 & 33 & $-4.356\times10^{-9}$ & $1.397\times10^{-10}$ & 1.0094 \\
14 & $-9.961\times10^{-4}$ & $7.324\times10^{-5}$ & 1.0535 & 34 & $-2.268\times10^{-9}$ & $6.985\times10^{-11}$ & 1.0088 \\
15 & $-5.347\times10^{-4}$ & $3.662\times10^{-5}$ & 1.0466 & 35 & $-1.187\times10^{-9}$ & $3.492\times10^{-11}$ & 1.0084 \\
16 & $-2.856\times10^{-4}$ & $1.831\times10^{-5}$ & 1.0409 & 36 & $-6.271\times10^{-10}$ & $1.746\times10^{-11}$ & 1.0079 \\
17 & $-1.520\times10^{-4}$ & $9.155\times10^{-6}$ & 1.0361 & 37 & $-3.124\times10^{-10}$ & $8.731\times10^{-12}$ & 1.0075 \\
18 & $-8.057\times10^{-5}$ & $4.578\times10^{-6}$ & 1.0322 & 38 & $-1.619\times10^{-10}$ & $4.366\times10^{-12}$ & 1.0071 \\
19 & $-4.257\times10^{-5}$ & $2.289\times10^{-6}$ & 1.0288 & 39 & $-8.226\times10^{-11}$ & $2.183\times10^{-12}$ & 1.0067 \\
20 & $-2.243\times10^{-5}$ & $1.144\times10^{-6}$ & 1.0260 & 40 & $-4.882\times10^{-11}$ & $1.091\times10^{-12}$ & 1.0064 \\
\bottomrule
\end{tabular}
\end{center}

\subsection{Time-dependent ramp}
\paragraph{Ramp-time scan (\code{run\_ramp}, smoothstep, cell 31).}
\begin{center}\small
\begin{tabular}{@{}rlllll@{}}
\toprule
$T$ & leak & wrong-branch weight & $\Delta Q(T)$ & \code{ΔQ\_adiabatic} & $\Delta Q(T)-$\code{ΔQ\_adiabatic} \\
\midrule
150 & $1.291\times10^{-7}$ & $9.0\times10^{-18}$ & $-0.889103$ & $-0.889108$ & $4.43\times10^{-6}$ \\
200 & $5.326\times10^{-8}$ & $2.2\times10^{-21}$ & $-0.889106$ & $-0.889108$ & $1.98\times10^{-6}$ \\
250 & $1.495\times10^{-8}$ & $1.2\times10^{-21}$ & $-0.889105$ & $-0.889108$ & $2.37\times10^{-6}$ \\
300 & $5.493\times10^{-9}$ & $1.2\times10^{-22}$ & $-0.889103$ & $-0.889108$ & $4.97\times10^{-6}$ \\
350 & $4.159\times10^{-9}$ & $3.0\times10^{-22}$ & $-0.889100$ & $-0.889108$ & $7.10\times10^{-6}$ \\
\bottomrule
\end{tabular}
\end{center}
The leak decreases as the ramp is slowed down; $\Delta Q(T)$ converges to $-0.88910$. The adiabatic value
$\Delta Q(\psi_+)=-0.889108$ is the same for every $T$ (the final Hamiltonian does not depend on $T$). The remaining
difference $\Delta Q(T)-\Delta Q(\psi_+)$ of 2--7$\times10^{-6}$ does not decrease with $T$ while the leak falls by a
factor of 30, so it is not a non-adiabatic effect; it is consistent with the accumulated time-step error of the
midpoint-rule TDVP ($dt=0.1$, more steps for larger $T$).

\paragraph{Checks and full densities.}
\begin{itemize}
\item Initial mode: $\lVert H(0)\psi_f\rVert=1.12\times10^{-12}$, $\Delta Q(0)=4.6\times10^{-15}$,
      $\lVert P_{\text{bulk}}\psi_f\rVert^2=1.8\times10^{-25}$.
\item Fermi level for the ramp: $\epsilon_F=-0.58831$ in the window $(-0.82787,-0.34875)$; $\max|E|=7.5749$.
\item $T=100$, TDVP bulk: overlap error $9.9\times10^{-9}$, completeness error $4.3\times10^{-9}$.
\item $\max|n_{\text{bulk}}^{\text{TDVP}}-n_{\text{bulk}}^{\text{McWeeny}}| = 1.29\times10^{-3}$.
\item Zero modes (cell 33): $T=300$: $\lVert C\psi_f(T)-\bar\psi_f(T)\rVert=9.5\times10^{-14}$,
      $|\langle\psi_+\vert\psi_1(T)\rangle|^2=0.9999999945$, $|\langle\psi_+\vert\psi_0(T)\rangle|^2=1.2\times10^{-22}$;
      $T=100$: $4.5\times10^{-14}$, $0.9999942664$, $3.7\times10^{-13}$. The same $T=100$ values are printed in cell 39.
      $\psi_f$ and $\bar\psi_f$ evolved separately stay charge-conjugate partners to $10^{-13}$, and the missing
      weight $1-|\langle\psi_+\vert\psi_1(T)\rangle|^2$ ($5.7\times10^{-6}$ for $T=100$, $5.5\times10^{-9}$ for $T=300$)
      is the non-adiabatic leak.
\item Method 1 vs method 2 over the $\ell$ sweep (cell 24): $\max_\ell|\Delta Q^{(1)}-\Delta Q^{(2)}|=5.6\times10^{-16}$.
\end{itemize}

\end{document}
